% SPDX-License-Identifier: GPL-3.0-or-later
% Copyright (C) 2026 Alessandro Greco
\chapter{前言}

每天都有数百万人使用 Tor，但人们谈论它时，往往首先想到的是它的用途——无论这些用途是真实存在的，还是出于想象：绕过审查、保护消息来源、所谓的 \enquote{暗网}。本书采取另一条路径：解释 Tor \emph{究竟如何工作}。我们从两个中继之间交换的字节开始，一章一章地向上追踪，直到设计取舍、攻击、现实案例以及仍未解决的问题。更重要的是，本书并不要求读者仅凭文字相信这些解释：技术部分所描述的每一种机制，都会在实验中构建并验证，实验运行在一张普通计算机即可承载的私有 Tor 网络上。

\section*{来源、创作方法与勘误邀请}

本书以我多年来撰写的学习笔记和研究资料为基础，最初的研究始于本科毕业论文。为了将这些资料整理成一部连贯的技术著作，我在自己的指导下使用人工智能工具，实质性地参与了讲解的扩展、论证、正文撰写，以及插图、习题、代码和实验的制作。人工智能并非只用于润色我写好的文字，其参与也不等于独立专家审查。

尽管已进行技术检查并由我本人审阅，书中仍可能存在错误、简化之处或过时的信息。欢迎读者通过项目仓库报告问题并提出修正。我选择通过项目官方渠道免费提供本书及相关源码，并采用 GNU GPL-3.0-or-later 许可证，以践行 \enquote{diffusione dei saperi senza fondare poteri}——传播知识，而不建立权力中心。包括商业再分发在内的使用条件由许可证规定。

\section*{本书适合谁}

本书面向计算机科学与安全专业的学生、工作中会接触 Tor 的专业人员、正在运行或希望运行中继的人，以及研究人员。读者应具备基本的网络知识（TCP/IP，以及 TLS 的基本思想），熟悉命令行；要完成实验，还需要掌握 Python 基础。本书所需的密码学知识会在 \cref{cap:05-crittografia} 中介绍，并在 \cref{cap:B-crittografia} 中进一步展开。

\section*{如何阅读本书}

全书分为五个部分。你可以从头到尾顺序阅读，也可以按照以下路径选择重点：
\begin{description}
  \item[历史与概念路径] 第一部分（第 1--4 章）和第五部分（第 29--31 章）：匿名究竟意味着什么，Tor 从何而来，它针对哪些对手提供保护，以及它正在走向何处。
  \item[协议路径] 第二部分（第 5--16 章），建议按顺序并结合实验阅读。这是本书的核心。
  \item[实践路径] 第三部分（第 17--25 章）：配置 tor、观察其流量、发布洋葱服务、运行中继、使用 Tor Browser、测量网络并理解相关法律问题。\cref{cap:19-rete-privata} 是所有实验的前置章节。
  \item[安全路径] 第四部分（第 26--28 章），建议先阅读 \cref{cap:04-modello-minaccia}：从研究成果、现实事件与运维失误三个角度分析攻击。
\end{description}

\section*{约定}

每一章都会先列出学习目标，并给出一行\emph{最后核验}信息，说明本章内容是针对哪个日期、软件版本和规范版本核对的。Tor 会不断演进，因此读者必须知道正文所描述的是哪一个版本。全书中，\emph{Tor} 指网络和项目，\emph{tor} 指 C 语言实现，\emph{Arti} 指 Rust 实现。规范引用写作 \spec{document}{section}；协议变更提案则以编号标识。正文采用的核心英文协议术语及其简体中文译法统一收录于 \cref{cap:A-glossario}。

为了不打断主线叙述，本书使用不同的框区分\emph{深入说明}、\emph{历史说明}和\emph{注意事项}。练习按星级标注难度：一星表示理解题，两星表示应用题，三星表示实验任务或小型研究项目；完整解答见 \cref{cap:D-soluzioni}。全书共有 20 个实验。在第二部分中，读者将用 Python 逐步构建一个小型 Tor 客户端（\texttt{torlab}）；代码与实验网络随书提供，并可在更新后重新运行自动化测试。所有这些资源——\texttt{torlab} 包、20 个实验、Docker 测试网络和验证工具——都可以在本书仓库中获得：\repolibro。

\section*{伦理说明}

本书的全部实验要么运行在实验室的私有 Tor 网络上，要么仅以普通客户端身份使用真实网络，不观察也不修改他人的流量。这既符合 Tor Research Safety Board 对研究人员的建议，也涉及法律边界（见 \cref{cap:25-legale}）。讨论攻击的章节始终采取分析和防御视角：解释过去发生过什么、如何防御，而不是说明如何针对具体个人实施攻击。

\section*{致谢}

本书建立在 Tor Project 及其社区长期工作的基础之上。开放规范、\texttt{tor} 与 Arti 的源代码、协议变更提案和公开文档，为后续每一项技术论断提供依据。正因为这些材料公开可得，读者才能回到原始来源核验，而不必仅仅相信作者的表述。同样需要感谢用于撰写本书和测试实验的自由软件项目及其维护者，包括 \LaTeX{}、Python、Docker、Chutney、Wireshark，以及正文中引用的许多其他项目。
